\subsection{Ray conditions at three and exact lower-norm repair}

\begin{statement}{A.124}{Proved}
{The ordinary permutation kernel excludes every deeper pure radical}
Let \(a\in\Q^\times\) have \(v_3(a)=1\), let
\(K_n=\Q(a^{1/3^n})\), and put
\(L_n=K_n(\mu_{3^n})\), with \(n\ge1\).
Let \(H_0/K_n\) be any finite abelian three-extension,
and define its actual normal compositum
\[
 F=L_n\prod_{\sigma\in G_\Q}\sigma H_0.
 \tag{A.124.1}
\]
For \(t>n\) and any cyclotomic \(R\) containing
\(\mu_{3^t}\), choose compatible roots and put
\(L_j=\Q(a^{1/3^j},\mu_{3^j})\).
Then
\[
                         FR\cap L_tR=L_nR.
 \tag{A.124.2}
\]
No restriction at three on the ramification of
\(H_0/K_n\) is needed.

In particular, retain the actual coefficients and
canonical representative \(\beta\) of A.101.
The simultaneous field \(S_t\) of A.105 can be used
with every three-primary ray modulus of \(K_n\)
supported away from two, including arbitrarily high
powers of the primes above three.
For the normal ray compositum (A.124.1), the exact maps are
\[
                  FR\cap S_t=L_nR,\qquad
                            S_t^{\rm ab/\Q}=R.
 \tag{A.124.3}
\]
Consequently the ordinary and same-lower-prime
relative Artin generators in A.99--A.100 retain
their complete boundary, prior-representative,
deeper-\(a\)-root, and compatible cyclotomic
splitting conditions with these ray moduli.

There is also an exact refinement of the relative
repair when the lower step has degree three.
Let \(\beta_{n-1}=N_{K_n/K_{n-1}}\beta\), and
apply the strengthened A.100 at precision \(3^r\),
\(r\ge2\), to get
\[
 N_{K_n/K_{n-1}}\beta'
              =\beta_{n-1}\eta^{3^r},
                  \qquad\eta\in\mathcal O_{K_{n-1}}^\times.
 \tag{A.124.4}
\]
Define the actual new representative
\[
                      \beta''=\beta'\eta^{-3^{r-1}}.
 \tag{A.124.5}
\]
Then
\[
 \begin{gathered}
 N_{K_n/K_{n-1}}\beta''=\beta_{n-1},\qquad
              \operatorname{div}(\beta'')=\operatorname{div}(\beta'),\\
 \beta''/\beta'\in(K_n^\times)^{3^{r-1}} .
 \end{gathered}
 \tag{A.124.6}
\]
Thus every lower representative in this existing
finite norm chain is retained exactly.
The ray modulus can be chosen so that \(\beta''/\beta\)
belongs to any prescribed neighborhood of one at
the unique place above three.
All the canonical two-adic valuations are retained.
The newly permitted primes and the valuation
divisibility conclusion of A.100 are unchanged.
\end{statement}

\begin{proof}
\emph{The actual restriction embedding.}
Put \(N=3^n\),
\[
 G=\operatorname{Gal}(L_n/\Q)=C_N\rtimes U_N,\quad
 U_N=\operatorname{Gal}(L_n/K_n),\quad
 A=\operatorname{Gal}(H_0/K_n),\quad
 V=\operatorname{Gal}(F/L_n).
 \tag{A.124.7}
\]
The group \(V\) is abelian and three-primary:
over \(L_n\), every conjugate of \(H_0\) gives
an abelian three-extension.
Let \(\widetilde G=\operatorname{Gal}(F/\Q)\)
and \(\widetilde U=\operatorname{Gal}(F/K_n)\).
Restriction gives
\(\rho:\widetilde U\to A\).
For \(v\in V\) and \(gU_N\in G/U_N\), choose a
lift \(\widetilde g\in\widetilde G\) and set
\[
  j(v)_{gU_N}
           =\rho(\widetilde g^{-1}v\widetilde g)
          \quad\text{in }A.
 \tag{A.124.8}
\]
This is independent of both choices.
Changing the lift by an element of \(V\) changes
the restricted element only by conjugation in
the abelian group \(A\).
Changing \(g\) by an element of \(U_N\) is
conjugation by an element of \(\widetilde U\),
which has the same property.
These facts also prove
\[
               j:V\lhook\joinrel\longrightarrow A^{G/U_N}
       \quad\text{is an equivariant homomorphism}
 \tag{A.124.9}
\]
for the ordinary permutation action.
The kernel fixes \(L_n\) and every conjugate of
\(H_0\), hence fixes \(F\) and is trivial.
The conjugation action on \(V\) factors through
\(G\), since \(V\) is abelian.
No linear disjointness of \(H_0\) and \(L_n\)
was used; the separate restrictions can simply
have images smaller than \(A\).

The translation \(C_N\) is transitive on \(G/U_N\).
Its fixed elements in the target of (A.124.9)
are therefore diagonal tuples.
Every such tuple is fixed by \(G\), in particular
by complex conjugation \(\epsilon\).
Consequently
\[
                             V^{C_N}\subset V^{\epsilon}.
 \tag{A.124.10}
\]
This is an ordinary permutation module statement.
There is no cyclotomic twist on its target.

\emph{Cyclotomic base change and the lift power.}
Restriction identifies
\[
 V_R=\operatorname{Gal}(FR/L_nR)
        \simeq\operatorname{Gal}(F/F\cap L_nR)
                          \subset V .
 \tag{A.124.11}
\]
This is the stable subgroup required after base
change, rather than a quotient of \(V\).
The full radical degree gives
\(\operatorname{Gal}(L_nR/R)=C_N\).
Its generator induces the original translation
on the subgroup in (A.124.11), so its fixed
elements remain \(\epsilon\)-fixed by (A.124.10).

Suppose \(L_{n+1}R\subset FR\).
Choose \(\sigma\in\operatorname{Gal}(FR/R)\)
lifting a generator of \(\operatorname{Gal}(L_nR/R)\).
Then \(\sigma^N\in V_R\), and it is fixed by
conjugation by \(\sigma\), hence by the translation.
It is therefore fixed by \(\epsilon\).
On the other hand its restriction to \(L_{n+1}R\)
is a nonzero element of the order-three kernel in
\[
 1\longrightarrow C_3\longrightarrow
           \operatorname{Gal}(L_{n+1}R/R)=C_{3N}
                \longrightarrow C_N\longrightarrow1.
 \tag{A.124.12}
\]
Indeed every lift of a generator of \(C_N\)
has order \(3N\) in this cyclic group.
Complex conjugation inverts that kernel:
it fixes the chosen real odd root of \(a\)
and inverts the roots of unity.
Restriction of \(\sigma^N\) would therefore
be both fixed and negated.
Its order is three, which makes this impossible.
Thus \(L_{n+1}R\) is not contained in \(FR\).
The intermediate fields of the cyclic extension
\(L_tR/L_nR\) form a chain, and every nontrivial
one contains \(L_{n+1}R\).
This proves (A.124.2).

\emph{The simultaneous boundary fields and actual ray maps.}
Use the coefficients of A.101, so \(a\) is
a two-adic unit.
If the ray modulus has no prime above two,
\(F/L_n\) is unramified there.
After base change to \(L_tR\) it remains unramified.
At two, \(L_tR/R\) splits completely: the odd
power map on \(\Q_2^\times\) gives the required
root of the unit \(a\).
The canonical two-adic divisor of the prior
\(\beta\) therefore gives exactly the repeated
valuation matrix of A.105.6.
The simultaneous \(\beta\)-root and \(b\)-root
lattice remains injective in those valuations,
including after its prime-to-three boundary
base change.
Every nonzero cyclic three-primary subextension
ramifies at some prime above two.
The unramified ray compositum can have no such
intersection; the prime-to-three part is disjoint
by degree.
Thus
\[
                   (FL_tR)\cap S_t=L_tR.
 \tag{A.124.13}
\]
Intersect this equality with \(FR\), and apply
(A.124.2), to obtain the first equality of
(A.124.3).
The second is the literal arithmetic group
calculation of A.101 and A.105, which does
not depend on the ray modulus.

These are precisely the maps used by the
Artin-character proofs.
For clarity, the permitted Frobenius elements
in \(\operatorname{Gal}(F/L_n)\) still form
the coset \(\sigma_0V_0\) of A.95.
Its quotient \(V/V_0\) comes from the abelian
intersection with the cyclotomic field and
has trivial conjugation action.
The same character-annihilator argument in
A.99 gives every ordinary split-prime target.
With \(J=\operatorname{Gal}(F/K_{n-1})\),
the same argument in A.100 gives differences
above the same lower prime.
Its extension map uses
\(H^2(J/V,\Q_3/\Z_3)=0\) and the actual
surjection \(H/V\to(J/V)^{\rm ab}\), unchanged
by allowing ray powers above three.
Chebotarev supplies those cosets with the
specified prior-field splitting, any compatible
larger cyclotomic condition, and the finite
exclusions already allowed in A.100.
This proves the asserted strengthening of
the actual ray Artin maps.

\emph{The exact lower norm and the three-adic congruence.}
Apply that relative ray proof with arbitrarily
high ray precision above three and with the
finite unit detectors of A.100.
It gives (A.124.4).
The lower field has index three in the top
field, so the norm of its scalar \(\eta\)
is \(\eta^3\).
Consequently
\[
 N_{K_n/K_{n-1}}(\eta^{-3^{r-1}})
                         =\eta^{-3^r},
 \tag{A.124.14}
\]
which cancels the factor in (A.124.4) exactly.
Since \(\eta\) is a global unit, it changes
no finite valuation.
Since \(r\ge2\), its exponent is a multiple
of \(3^{r-1}\), giving the final assertion
of (A.124.6).
Transitivity of field norms retains every
lower member of the existing norm chain
and its exact rational norm.
The canonical two-adic valuations are among
those unchanged valuations.

At the unique three-adic completion of
\(K_{n-1}\), there are no nontrivial
three-power roots of unity: its degree
over \(\Q_3\) is odd, whereas \(\mu_3\)
has degree two.
The power map \(z\mapsto z^{3^r}\) on
its compact unit group is therefore
injective and is a homeomorphism onto
its image.
The ray condition can make
\(\beta'/\beta\) arbitrarily close to one.
Its relative norm \(\eta^{3^r}\) then
tends to one, so the unique local root
\(\eta\) tends to one.
Formula (A.124.5) now makes
\(\beta''/\beta\) arbitrarily close to
one as well.
This proves the claimed three-adic
congruence without asserting an
uncontrolled global choice of a root.
\end{proof}

The exact multiplier entering A.122 is
\(\beta''/\beta\in T_n(\Q)\).
The new argument permits arbitrary ray
precision at three while retaining the
prescribed deeper root fields and exact
lower norm.
Its old local quadratic values must still
be evaluated through (A.122.3); the ray
argument alone does not make every one
of those values zero.
Likewise it repairs one existing norm
representative and does not manufacture
a next preimage of an arbitrarily fixed
representative.

The retained exact certificate
\path{code/ray_three_permutation_certificate.py}
checks the ordinary permutation actions,
their diagonal translation invariants,
the larger cyclotomic actions, and the
lift-power contradiction for finite
three-primary coefficient groups.
